% answers_short.tex —— 2005 年·普通高等学校招生考试·湖北卷·文科 · 简短答案速查
% 由 tools/make_short_answers.py 自 answers.tex 的 \gksolhead/\gkanswer 提取，
% 并与 2005/2005_数学真题_湖北文科_答案要点.md 逐题比对后生成，勿手改（改源 MD 后重跑）。

\begin{gknotice}
\textbf{使用说明}：本节只给最终答案便于快速核对，完整推导见解析版。答案为编写者独立推导，\textbf{非官方标准答案}。
\end{gknotice}

\gkpart{选择题}
\begin{center}\begin{tabular}{ccccccccccccc}
\toprule
题号 & 1 & 2 & 3 & 4 & 5 & 6 & 7 & 8 & 9 & 10 & 11 & 12 \\
\midrule
答案 & B & B & C & D & C & A & B & A & D & C & D & D \\
\bottomrule
\end{tabular}\end{center}

\gkpart{填空题}
\begin{itemize}[leftmargin=2.6em]
  \item[\textbf{13.}] $[2, 3) \cup (3, 4)$（或 $\{x \mid 2 \leqslant x < 4 \text{ \text{且} } x \neq 3\}$）。
  \item[\textbf{14.}] $38$。
  \item[\textbf{15.}] $2\pi - \dfrac{1}{2}$。
  \item[\textbf{16.}] $500$。
\end{itemize}

\gkpart{解答题}
\begin{itemize}[leftmargin=2.6em]
  \item[\textbf{17.}] $[5, +\infty)$。
  \item[\textbf{18.}] $6\sqrt{2} + 8\sqrt{3}$。
  \item[\textbf{19.}] (1) $a_n = 4n - 2$；$b_n = 2 \times \left(\dfrac{1}{4}\right)^{n-1} = \dfrac{1}{2^{2n-3}}$；(2) $T_n = \dfrac{(6n - 5) \cdot 4^n + 5}{9}$
  \item[\textbf{20.}] (1) $BF = 2\sqrt{6}$；(2) $\dfrac{4\sqrt{33}}{11}$
  \item[\textbf{21.}] (1) 不需要换灯泡的概率为 $p_1^5$；更换 2 只灯泡的概率为 $10(1 - p_1)^2 p_1^3$；(2) $(1 - p_1)^2 + p_1 - p_2$；(3) $0.24$
  \item[\textbf{22.}] (1) $\lambda$ 的取值范围是 $(12, +\infty)$；直线 $AB$ 的方程为 $x + y - 4 = 0$；(2) 存在，四点 $A, B, C, D$ 对任意 $\lambda \in (12, +\infty)$ 恒共圆，证明见解析
\end{itemize}
